\documentclass[11pt,a4paper]{article}

\usepackage[margin=2.4cm]{geometry}
\usepackage{amsmath,amssymb}
\usepackage{booktabs}
\usepackage{xcolor}
\usepackage{tcolorbox}
\usepackage{fancyhdr}
\usepackage{enumitem}
\usepackage{listings}
\usepackage[hidelinks]{hyperref}

\tcbuselibrary{skins,breakable}

\definecolor{keyblue}{HTML}{1F4E79}
\definecolor{warnamber}{HTML}{8A6100}
\definecolor{okgreen}{HTML}{1E6B3A}
\definecolor{softgrey}{HTML}{F4F5F7}
\definecolor{rulegrey}{HTML}{C9CDD3}
\definecolor{codebg}{HTML}{F7F7F9}

\pagestyle{fancy}
\fancyhf{}
\lhead{\small CS F407 --- Logical Reasoning for Planning}
\rhead{\small Personal study notes}
\cfoot{\small\thepage}
\renewcommand{\headrulewidth}{0.4pt}

\newtcolorbox{keyidea}[1]{
  breakable, colback=softgrey, colframe=keyblue,
  boxrule=0.9pt, left=8pt, right=8pt, top=6pt, bottom=6pt,
  title={\bfseries #1}, fonttitle=\small\bfseries
}

\newtcolorbox{warnbox}[1]{
  breakable, colback=softgrey, colframe=warnamber,
  boxrule=0.9pt, left=8pt, right=8pt, top=6pt, bottom=6pt,
  title={\bfseries #1}, fonttitle=\small\bfseries
}

\newtcolorbox{okbox}[1]{
  breakable, colback=softgrey, colframe=okgreen,
  boxrule=0.9pt, left=8pt, right=8pt, top=6pt, bottom=6pt,
  title={\bfseries #1}, fonttitle=\small\bfseries
}

\lstset{
  basicstyle=\ttfamily\small,
  backgroundcolor=\color{codebg},
  frame=single, rulecolor=\color{rulegrey},
  breaklines=true, showstringspaces=false,
  xleftmargin=0pt, framexleftmargin=3pt,
  aboveskip=8pt, belowskip=8pt,
}

\setlist[itemize]{itemsep=2pt, topsep=4pt}
\setlist[enumerate]{itemsep=2pt, topsep=4pt}

\setlength{\emergencystretch}{3em}

\newcommand{\ent}{\models}

\title{\vspace{-1.5cm}\textbf{Logic, Search and Planning:\\An End-to-End Explainer}}
\author{The Warehouse Robot --- CS F407 Artificial Intelligence}
\date{Study notes, not for submission}

\begin{document}
\maketitle

\begin{keyidea}{How to read this}
This is the whole laboratory explained from first principles: what a planning
problem is, what it means for an action to be \emph{applicable}, how logic and
search fit together, what the code actually does, and --- the part that matters
--- why one of the bugs in this laboratory was invisible to \emph{every} check
that could be written in Python.

Every number quoted here was measured from the programs in
\texttt{deliverables/}. Nothing is illustrative.

If you read only one section, read \S7. It is the part of this laboratory that
teaches something a textbook will not.
\end{keyidea}

\tableofcontents
\vspace{6pt}
\hrule

% =====================================================================
\section{The one-sentence version}

\begin{keyidea}{The whole laboratory}
\[
\boxed{\text{Logic} + \text{Search} = \text{Planning}}
\]
\emph{Logic} decides what is \textbf{possible}: given the state of the world,
which actions can be performed, and what does the world look like afterwards?
\emph{Search} decides what to \textbf{try}: of all the sequences of possible
actions, which one do I explore next?

A planner is a graph search in which the graph is not given to you --- it is
\emph{computed}, one node at a time, by a logical test.
\end{keyidea}

That last sentence is the difference from the previous module. In the search
laboratory you were handed a grid: the neighbours of a cell were simply the
adjacent cells. Here nobody hands you the graph. You describe the world with
propositions and actions, and the edges of the graph come into existence
whenever an action's preconditions happen to be satisfied.

This has a consequence that runs through everything below: \textbf{you can now
get the graph itself wrong}, and no amount of searching it correctly will tell
you so.

% =====================================================================
\section{States: the world as a set of true things}

\subsection{The representation}

A \textbf{proposition} is an atomic statement that is either true or false:
$\mathrm{At}(Robot, A)$, $\mathrm{Holding}(Package)$. A \textbf{state} is the
set of propositions that are \emph{true}:
\[
  S_0 = \{\, \mathrm{At}(Robot,A),\ \mathrm{At}(Package,A) \,\}.
\]

In code, a state is a Python \texttt{frozenset} of strings.

\begin{warnbox}{The closed-world assumption --- the most important line in the file}
Anything \emph{not} in the set is \textbf{false}. Not ``unknown''---false.

$S_0$ does not contain $\mathrm{Holding}(Package)$, and that is precisely how
the program knows the robot's hands are empty. It does not contain
$\mathrm{At}(Robot,B)$, and that is how it knows the robot is not at B.

This is why a state of 2 propositions can pin down a world of 7 facts. It is
also why, later, Prolog answering \texttt{false} to \texttt{can\_move(a,c)}
means ``I cannot derive it'', which the system chooses to treat as ``it is
false''. Same assumption, same trade.
\end{warnbox}

\subsection{Why \texttt{frozenset} and not \texttt{set}}

Because states are used as \emph{keys}: they go into a \texttt{visited} set and
a \texttt{parent} dictionary. Python sets are mutable, therefore unhashable,
therefore cannot be dictionary keys. The first LLM-generated version used plain
sets and died with
\begin{lstlisting}
TypeError: unhashable type: 'set'
\end{lstlisting}
the moment duplicate detection was added. Immutability is not a nicety here: a
state that can be mutated after being recorded as visited would silently corrupt
the search.

\subsection{A modelling choice worth understanding}

Where does the package go when the robot picks it up? The representation says:
\emph{nowhere}. $\mathrm{At}(Package,A)$ is deleted and
$\mathrm{Holding}(Package)$ is added. While it is held, the package has
\textbf{no location proposition at all}.

That sounds like a gap. It is the opposite --- it is what makes \texttt{Move}
correct for free:

\begin{itemize}
  \item If the package is on the floor and the robot walks away, the package's
        $\mathrm{At}$ proposition is untouched, so it stays behind. Correct.
  \item If the package is held, there is no $\mathrm{At}$ proposition to
        update, so it travels with the robot. Also correct.
\end{itemize}

One \texttt{Move} definition, both behaviours, no special cases. The
alternative --- $\mathrm{At}(Package, Robot)$ --- would need \texttt{Move} to
know about the package, and would admit states where the package is both at A
and in the robot's hand.

\begin{okbox}{Invariant}
$\mathrm{Holding}(Package)$ and $\mathrm{At}(Package,x)$ are never both true and
never both false. The package is always in exactly one of: a place, or the
robot's hand.
\end{okbox}

% =====================================================================
\section{Actions: preconditions and effects}

An action has four sets and a name:

\begin{center}
\begin{tabular}{@{}ll@{}}
\toprule
$\mathrm{Pre}^{+}$ & must be \textbf{present} in the state \\
$\mathrm{Pre}^{-}$ & must be \textbf{absent} from the state \\
$\mathrm{Add}$     & made \textbf{true} by the action \\
$\mathrm{Del}$     & made \textbf{false} by the action \\
\bottomrule
\end{tabular}
\end{center}

\subsection{Applicability: the logical step}

\begin{keyidea}{The entailment check}
\[
  S \ent \mathrm{Pre}(a)
  \qquad\Longleftrightarrow\qquad
  \mathrm{Pre}^{+}(a) \subseteq S
  \ \text{ and }\
  \mathrm{Pre}^{-}(a) \cap S = \varnothing .
\]
In Python, one line:
\begin{lstlisting}
return self.pos_pre <= state and not (self.neg_pre & state)
\end{lstlisting}
\end{keyidea}

This is what ``logical reasoning'' means in this laboratory. It is not a theorem
prover. It is a small \emph{model-checking} step --- is this formula true in
this model? --- performed once per action per state. Solving the warehouse
performs it \textbf{70 times}.

The reason it is so cheap is that the actions are \emph{ground}: all their
variables have already been filled in, so $\mathrm{Move}(A,B)$ and
$\mathrm{Move}(B,C)$ are separate objects. Real planners keep the schemas and
unify at search time, which is more expressive and much more expensive.

\subsection{Effects: the transition}

\[
  S' = \mathrm{Apply}(S, a) = \bigl(S \setminus \mathrm{Del}(a)\bigr) \cup \mathrm{Add}(a).
\]

\begin{warnbox}{Delete first, \emph{then} add --- and why the order is specified}
If a proposition appears in both $\mathrm{Add}$ and $\mathrm{Del}$, the order
decides its fate:
\begin{itemize}
  \item delete-then-add: it \textbf{survives};
  \item add-then-delete: it is \textbf{lost}.
\end{itemize}
No action in this domain has such an overlap, so both orders give identical
results \emph{here}. That is exactly why it must be fixed in the specification
rather than discovered later: an underspecified detail that happens not to
matter today is a bug waiting for the domain to grow. The prompt in
\texttt{prompts.md} states the order explicitly so the generated code could not
choose for itself, and Test~G pins it down.
\end{warnbox}

\subsection{The warehouse actions, in full}

\begin{center}
\small
\begin{tabular}{@{}lllll@{}}
\toprule
\textbf{Action} & $\mathrm{Pre}^{+}$ & $\mathrm{Pre}^{-}$ & $\mathrm{Add}$ & $\mathrm{Del}$ \\
\midrule
$\mathrm{Move}(x,y)$ & $\mathrm{At}(R,x)$ & --- & $\mathrm{At}(R,y)$ & $\mathrm{At}(R,x)$ \\
\emph{connected $x,y$ only} & & & & \\[4pt]
$\mathrm{PickUp}(P,x)$ & $\mathrm{At}(R,x)$, $\mathrm{At}(P,x)$ & $\mathrm{Holding}(P)$ & $\mathrm{Holding}(P)$ & $\mathrm{At}(P,x)$ \\[4pt]
$\mathrm{Drop}(P,x)$ & $\mathrm{At}(R,x)$, $\mathrm{Holding}(P)$ & --- & $\mathrm{At}(P,x)$ & $\mathrm{Holding}(P)$ \\
\bottomrule
\end{tabular}
\end{center}

Note the italicised line. \textbf{The connectivity of the warehouse is not
inside any action --- it decides which actions exist at all.} With
\[
  A \relbar\joinrel\relbar B \relbar\joinrel\relbar C, \qquad A \not\!\relbar\joinrel\relbar C,
\]
there are four \texttt{Move} actions, not six. Remember this sentence; \S7 is
about what happens when it is forgotten.

$\mathrm{PickUp}$'s negative precondition stops the robot picking up a package
it is already holding. It is the only $\mathrm{Pre}^{-}$ in the domain, and it
exists partly so that there \emph{is} one to test.

% =====================================================================
\section{The goal: entailed, not equalled}

\begin{keyidea}{$G \subseteq S$, never $S = G$}
\[
  G = \{\, \mathrm{At}(Package, C) \,\}, \qquad
  S_4 = \{\, \mathrm{At}(Package,C),\ \mathrm{At}(Robot,C) \,\}.
\]
$S_4 \supseteq G$, so $S_4 \ent G$: \textbf{goal achieved}. \\
$S_4 \neq G$: equality would say \textbf{failure}.
\end{keyidea}

The goal is a \emph{partial description of the world}. It constrains the package
and says nothing at all about the robot. Demanding equality would demand that
the final state contain the goal proposition \emph{and nothing else} --- which
is unachievable, because the robot is always somewhere.

This is the single most predictable place for generated code to go wrong, which
is why it was written down in Task~0 \emph{before} any prompting, and stated
explicitly in the prompt. It still had to be fixed once (Defect~0), and because
it had been specified, the fix was one line.

\subsection{Test on pop, not on generation}

The planner checks the goal when a state comes \emph{off} the frontier, not
when it is put on. For BFS with unit-cost actions both are correct, but testing
on pop means \texttt{states expanded} counts the same thing for every search
strategy, so BFS and DFS numbers can be compared honestly. (In the previous
module this mattered more: for A$^\ast$, goal-test-on-pop is what makes the
algorithm \emph{optimal}.)

% =====================================================================
\section{Search: walking a graph nobody built}

\subsection{The implicit graph}

The action definitions define a graph without ever constructing it:
\begin{itemize}
  \item \textbf{nodes} = states;
  \item \textbf{edge} $S \xrightarrow{\;a\;} S'$ exists exactly when
        $S \ent \mathrm{Pre}(a)$ and $S' = \mathrm{Apply}(S,a)$.
\end{itemize}

For the warehouse, exhaustively enumerated in \texttt{results.txt}~\S1, the
reachable graph has \textbf{12 nodes}. In general it has up to $2^{n}$ for $n$
propositions --- which is why the graph is generated lazily, on demand: the
applicability test \emph{is} the successor function.

\subsection{BFS, and what the visited set is really for}

\begin{lstlisting}
frontier = deque([start]);  visited = {start};  parent = {start: (None, None)}
while frontier:
    state = frontier.popleft()              # FIFO -> shortest plans first
    if problem.satisfies_goal(state): return reconstruct(parent, state)
    for action in problem.actions:
        if not action.applicable(state): continue      # <- the logic
        successor = action.apply(state)                # <- the logic
        if successor in visited: continue              # <- termination
        visited.add(successor); parent[successor] = (state, action)
        frontier.append(successor)
\end{lstlisting}

BFS returns a \textbf{shortest} plan because it explores states in order of
increasing plan length and every action costs 1 --- the same argument as in the
search module.

\begin{warnbox}{The \texttt{visited} set is not an optimisation}
It is a \textbf{termination requirement}, and the difference is not academic.

$\mathrm{Move}(A,B)$ and $\mathrm{Move}(B,A)$ form a 2-cycle. Without duplicate
detection, the frontier refills faster than it drains. On a \emph{solvable}
problem this does not matter --- BFS reaches the goal before it notices. On an
\emph{unsolvable} problem it never runs out of states, so it never reports
failure.

\begin{center}
\small
\begin{tabular}{@{}lll@{}}
\toprule
& \textbf{with \texttt{visited}} & \textbf{without} \\
\midrule
solvable problem   & 4-action plan, 8 expansions & 4-action plan, 25 expansions \\
unsolvable problem & \textbf{``no plan'', 3 expansions} & \textbf{still running at 50{,}000} \\
\bottomrule
\end{tabular}
\end{center}

Read the first row again: on the problem the handout calls Test~A, the broken
version and the correct version give \emph{the same right answer}. The bug is
only visible on Test~B. This is the entire reason the handout insists you test
the impossible case.
\end{warnbox}

\subsection{What changes if you swap the search, and what does not}

From \texttt{results.txt}~\S3, with the logic untouched:

\begin{center}
\small
\begin{tabular}{@{}llrrrr@{}}
\toprule
\textbf{Problem} & \textbf{Algo} & \textbf{Plan} & \textbf{Expanded} & \textbf{Generated} & \textbf{Peak frontier} \\
\midrule
warehouse            & BFS & \textbf{4} & 8  & 16  & 3 \\
warehouse            & DFS & 4          & 7  & 13  & 2 \\
corridor, 6 locations & BFS & \textbf{7} & 27 & 58  & 8 \\
corridor, 6 locations & DFS & 7          & 13 & 28  & 5 \\
corridor, 8 locations & BFS & \textbf{9} & 47 & 101 & 11 \\
corridor, 8 locations & DFS & 9          & 17 & 38  & 7 \\
\bottomrule
\end{tabular}
\end{center}

Both return \textbf{valid} plans, because both use the same applicability test.
Only BFS returns a \textbf{shortest} one. DFS matches it here only because a
corridor offers almost no wrong turns to take --- give it a branching map and it
will wander. DFS touches fewer states, which is the usual trade: it commits
early and is rewarded when the commitment happens to be right.

\begin{okbox}{The division of labour, made visible}
Changing the search changed \emph{which} plan was found and how much work it
took. It did not and could not change \emph{which plans are legal}. Logic fixes
the set of legal plans; search picks one out of it.
\end{okbox}

% =====================================================================
\section{Putting it together: the plan by hand}

\subsection{The trap in the handout}

The handout suggests you ``might need to reason about''
\[
  \mathrm{Move}(A,B),\quad \mathrm{PickUp}(Package,B),\quad
  \mathrm{Move}(B,C),\quad \mathrm{Drop}(Package,C).
\]

\textbf{That is not a valid plan}, and finding out why is the most useful five
minutes in the laboratory.

\begin{center}
\small
\begin{tabular}{@{}clll@{}}
\toprule
\textbf{Step} & \textbf{Action} & \textbf{State before} & \textbf{Applicable?} \\
\midrule
1 & $\mathrm{Move}(A,B)$        & $\{\mathrm{At}(R,A), \mathrm{At}(P,A)\}$ & yes \\
2 & $\mathrm{PickUp}(Package,B)$ & $\{\mathrm{At}(R,B), \mathrm{At}(P,A)\}$ & \textbf{no} --- $\mathrm{At}(P,B)$ is false \\
\bottomrule
\end{tabular}
\end{center}

The robot walked to B \emph{without the package}, because \texttt{Move} does not
move the package. The sequence names all the right actions in a sensible order
and is wrong at the second step.

\begin{keyidea}{The lesson, before an LLM is anywhere near the problem}
``Looks reasonable'' is a judgement about a \emph{list of action names}.
``Valid'' is a property of the list \emph{together with the states it passes
through}. The only way to tell them apart is to carry the state along and check
each precondition where the action is actually used.
\end{keyidea}

\subsection{The plan that works}

\begin{center}
\small
\begin{tabular}{@{}cll@{}}
\toprule
\textbf{State} & \textbf{Facts} & \textbf{Why the action was applicable} \\
\midrule
$S_0$ & $\mathrm{At}(R,A)$, $\mathrm{At}(P,A)$ & (given) \\
$S_1$ & $\mathrm{At}(R,A)$, $\mathrm{Holding}(P)$ & $\mathrm{PickUp}(P,A)$: $\mathrm{At}(R,A)$\checkmark\ $\mathrm{At}(P,A)$\checkmark\ $\mathrm{Holding}$ absent\checkmark \\
$S_2$ & $\mathrm{At}(R,B)$, $\mathrm{Holding}(P)$ & $\mathrm{Move}(A,B)$: $\mathrm{At}(R,A)$\checkmark \\
$S_3$ & $\mathrm{At}(R,C)$, $\mathrm{Holding}(P)$ & $\mathrm{Move}(B,C)$: $\mathrm{At}(R,B)$\checkmark \\
$S_4$ & $\mathrm{At}(R,C)$, $\mathrm{At}(P,C)$ & $\mathrm{Drop}(P,C)$: $\mathrm{At}(R,C)$\checkmark\ $\mathrm{Holding}(P)$\checkmark \\
\bottomrule
\end{tabular}
\end{center}

$S_4 \ent G$. \textbf{Four actions, and four is optimal:} one $\mathrm{PickUp}$
and one $\mathrm{Drop}$ are unavoidable (only $\mathrm{Drop}$ adds
$\mathrm{At}(P,C)$, and only $\mathrm{PickUp}$ adds the $\mathrm{Holding}$ it
requires), and two $\mathrm{Move}$s are unavoidable because A and C are not
adjacent. $2 + 2 = 4$.

Hold on to the number \textbf{4}. It is what catches the bug in the next
section.

% =====================================================================
\section{The bug that nothing inside Python could find}

\begin{keyidea}{If you remember one thing}
Three defects reached working code. They are worth ranking not by how bad they
look, but by \textbf{where the evidence that would catch them lives}.
\end{keyidea}

\subsection{Defect 0 --- the goal test written as \texttt{state == G}}

\textbf{Symptom:} ``No plan found'' on the obviously solvable problem.
\textbf{Caught by:} Test~A, instantly. \textbf{Evidence lives:} inside the code.

Loud, immediate, harmless. Predicted in writing before any prompt was issued.

\subsection{Defect 1 --- the missing \texttt{visited} set}

\textbf{Symptom:} the program hangs --- but only on the \emph{unsolvable}
problem. \textbf{Caught by:} Test~B. \textbf{Evidence lives:} inside the code,
but only if you run the right input.

Quieter. Test~A passes. It is worth noting that the LLM diagnosed the cause
correctly once told the symptom, and then described the fix as ``an efficiency
improvement'' --- a plausible sentence about a termination bug.

\subsection{Defect 2 --- \texttt{Move} over every pair of locations}

The generated code instantiated \texttt{Move} for all six ordered pairs of
locations instead of the four connected ones. So $\mathrm{Move}(A,C)$ existed.

\begin{center}
\small
\begin{tabular}{@{}lll@{}}
\toprule
& \textbf{correct domain} & \textbf{defective domain} \\
\midrule
ground actions            & 10 & 12 \\
plan                      & \texttt{PickUp(P,A)} & \texttt{PickUp(P,A)} \\
                          & \texttt{Move(A,B)}   & \textbf{\texttt{Move(A,C)}} \\
                          & \texttt{Move(B,C)}   & \texttt{Drop(P,C)} \\
                          & \texttt{Drop(P,C)}   & \\
length                    & 4 & \textbf{3} \\
passes \texttt{validate\_plan}? & yes & \textbf{yes} \\
shortest, per BFS?        & yes & \textbf{yes} \\
LLM asked to justify it?  & --- & \textbf{confirmed it, correctly} \\
possible in the warehouse? & yes & \textbf{no} \\
\bottomrule
\end{tabular}
\end{center}

\begin{warnbox}{Why every internal check passed --- and had to}
Preconditions, effects, the goal test, the independent re-execution and the
optimality proof are \emph{all evaluated relative to the action set}. The action
set was the thing that was wrong.

Re-executing the plan asks: \emph{is this plan consistent with the model?} Yes,
perfectly. Nobody was asking: \emph{is the model consistent with the
warehouse?}

There is no test you can write in Python, against this model, that rejects this
plan. The search is exhaustively, provably correct over a state space that was
described wrongly. Its correctness is not in question --- its \emph{subject} is.
\end{warnbox}

\textbf{How it was actually found:} by reading the output and noticing
\textbf{three} actions where the hand-built plan needed \textbf{four}. Not by a
test. By a \emph{number} that disagreed with something worked out beforehand ---
which is the entire practical value of doing Task~1 by hand before touching the
computer.

\subsection{The pattern}

\begin{center}
\small
\begin{tabular}{@{}llll@{}}
\toprule
& \textbf{Defect 0} & \textbf{Defect 1} & \textbf{Defect 2} \\
\midrule
kind of error       & logical      & algorithmic & \textbf{modelling} \\
symptom             & wrong answer & hang        & \textbf{none} \\
visible from inside? & yes         & yes         & \textbf{no} \\
found by            & Test A       & Test B      & a human, comparing with Task 1 \\
LLM's explanation   & silent       & ``an optimisation'' & \textbf{confidently confirmed it} \\
\bottomrule
\end{tabular}
\end{center}

\begin{keyidea}{The finding}
The LLM was reliable on everything whose correctness is judged \textbf{from
inside the code}, and unreliable on the one thing whose correctness is judged
\textbf{against the world} --- and it was \emph{most} confident precisely there.

The asymmetry is the finding. Not ``LLMs are unreliable''.
\end{keyidea}

% =====================================================================
\section{Why a generated explanation proves nothing}

Task~5 asks: which do you trust, the LLM's explanation, or the independently
executed state transitions? The expected answer is the transitions, and it is
right. This run produced a sharper version.

When the LLM was asked to justify the three-action plan, it produced a
step-by-step precondition check. \textbf{Every line of it was true.}
$\mathrm{PickUp}(P,A)$ needs $\mathrm{At}(R,A)$ and $\mathrm{At}(P,A)$: both
hold. $\mathrm{Move}(A,C)$ needs $\mathrm{At}(R,A)$: it holds.
$\mathrm{Drop}(P,C)$ needs $\mathrm{At}(R,C)$ and $\mathrm{Holding}(P)$: both
hold. The final state contains $\mathrm{At}(P,C)$.

The independent Python re-execution agreed. BFS additionally certified the plan
as \emph{shortest}. \textbf{Three confirmations, every one of them correct,
every one of them useless.}

\begin{keyidea}{An ordering of trust}
\begin{enumerate}
  \item \textbf{An independent check against an independent source of
        knowledge.} Prolog's \texttt{valid\_route(a,[move(a,c)])} $\to$
        \texttt{false}. The only thing that caught Defect 2.
  \item \textbf{Independently executed state transitions.} Catches every error
        \emph{within} the model --- mis-applied preconditions, mis-sequenced
        effects, a plan whose steps do not join up. Cannot catch an error
        \emph{in} the model.
  \item \textbf{A generated explanation.} Evidence of nothing. It is text
        conditioned on the plan, produced by the process that produced the plan,
        and it will be just as fluent when the plan is wrong.
\end{enumerate}
\end{keyidea}

The handout's own phrasing is the thing to memorise, with one clause added:

\begin{center}
\fbox{\parbox{0.86\textwidth}{\centering
A generated explanation is not the same as an independent verification\\
--- and an independent \emph{execution} is only as good as the model it executes.}}
\end{center}

% =====================================================================
\section{Prolog: an independent judge}

\subsection{What Prolog is, in three lines}

A Prolog program is \textbf{facts}, \textbf{rules} and \textbf{queries}.

\begin{lstlisting}
connected(a,b).                      % a fact
can_move(X,Y) :- connected(X,Y).     % a rule:  Connected(X,Y) -> CanMove(X,Y)
?- can_move(a,b).                    % a query: does this follow?
true.
\end{lstlisting}

\texttt{head :- body} is the implication $\mathit{body} \rightarrow
\mathit{head}$, written backwards. A query asks whether the goal is
\emph{derivable} from the facts and rules.

\subsection{How the answers are computed (SLD resolution)}

Prolog reasons \textbf{backwards}, from the goal:

\begin{center}
\texttt{can\_move(a,b)}
$\;\xrightarrow{\text{rule, } X\mapsto a,\, Y\mapsto b}\;$
\texttt{connected(a,b)}
$\;\xrightarrow{\text{fact}}\;$
$\square$ \quad(empty goal: success)
\end{center}

Each step replaces the current goal with the body of a matching clause, after
\textbf{unification} fills in the variables. Reach the empty goal and the query
succeeds. Run out of clauses and it backtracks.

For Task~8's chain, forward reasoning proves the conclusion:
\[
  \underbrace{\mathrm{WetRoad}}_{\text{fact}}
  \;\Longrightarrow\;
  \mathrm{Slippery}
  \;\Longrightarrow\;
  \mathrm{ReduceSpeed}
\]
by two applications of \emph{modus ponens}. Prolog reduces the \emph{goal}
\texttt{reduce\_speed} to \texttt{slippery} to \texttt{wet\_road}, and finds a
fact. \textbf{Same theorem, opposite direction of travel.}

\subsection{Why \texttt{can\_move(a,c)} fails, and why that is useful}

There is no fact \texttt{connected(a,c)}, and \texttt{can\_move/2} has one
clause whose body is a single \texttt{connected/2} goal --- it is \textbf{not
recursive, therefore not transitive}. No derivation exists, so Prolog answers
\texttt{false}.

\begin{warnbox}{``False'' means ``not derivable''}
This is \textbf{negation as failure}, the closed-world assumption again.
Classical logic would say \emph{unknown}; Prolog says \emph{false}.

Two consequences worth separating:
\begin{itemize}
  \item It is what makes Prolog usable as a verifier: an action the knowledge
        base cannot justify is treated as \emph{forbidden}. A verifier that
        answered ``unknown'' would verify nothing.
  \item It is also why Prolog is not classical logic. Its \emph{completion}
        semantics reads the clauses for \texttt{can\_move/2} as an
        if-and-only-if: $\mathrm{CanMove}(X,Y) \leftrightarrow
        \mathrm{Connected}(X,Y)$. The single implication alone would leave
        $\mathrm{CanMove}(a,c)$ undetermined.
\end{itemize}
\end{warnbox}

\subsection{Adjacent is not reachable}

\begin{lstlisting}
?- can_move(a,c).     false.      % a and c are not adjacent ...
?- reachable(a,c).    true.       % ... but c IS reachable from a, via b
\end{lstlisting}

\texttt{reachable/2} is the transitive closure and needs \emph{recursion}. These
are different questions --- and \textbf{Defect 2 is exactly the error of
treating the second as though it were the first.}

\subsection{The loop closed}

\texttt{run\_prolog.py} translates the Python plan into a Prolog goal and asks a
knowledge base that has never seen the planner:

\begin{lstlisting}
the corrected planner
    Python plan            : PickUp(Package,A) -> Move(A,B) -> Move(B,C) -> Drop(Package,C)
    Python self-validation : VALID
    translated Prolog goal : valid_route(a, [move(a,b),move(b,c)]).
    Prolog answer          : true.        >>> AGREEMENT.

the defective planner (Move over every pair of locations)
    Python plan            : PickUp(Package,A) -> Move(A,C) -> Drop(Package,C)
    Python self-validation : VALID
    translated Prolog goal : valid_route(a, [move(a,c)]).
    Prolog answer          : false.       >>> DISAGREEMENT.
            valid_move(a,c) -> FALSE  <-- unsupported
\end{lstlisting}

The second block is the whole laboratory in six lines. Python says the plan is
consistent with its own action model; Prolog says the action model is not
consistent with the warehouse. \textbf{The verifier is right.}

\begin{okbox}{Why this works, and when it would not}
The two sides share \textbf{no code and no data}. The connectivity of the
warehouse is written down twice --- as \texttt{CONNECTED} in Python, as
\texttt{connected/2} facts in Prolog --- and the verifier earns its keep exactly
where the two copies disagree.

Had the Prolog facts been \emph{generated from} the Python action list, the
check would have been theatre: it would agree by construction, including when
both were wrong. \textbf{The independence is the mechanism, not a detail of it.}
\end{okbox}

\subsection{The general architecture}

\begin{center}
\fbox{\ Generate\ (fast, fallible)\ $\longrightarrow$\ Verify\ (narrow, trustworthy)\ }
\end{center}

The verifier does not need to be able to \emph{write} a plan --- it only needs
to be able to \emph{reject} a bad one, which is a far easier job. That
asymmetry is why the pattern scales: type checkers, proof assistants,
model checkers and unit tests are all the same shape.

% =====================================================================
\section{Answers you should be able to give}

\begin{enumerate}\itemsep6pt

\item \textbf{Is $\mathrm{PickUp}(Package,A)$ applicable in $I$?}
Yes. $\mathrm{At}(R,A)$ and $\mathrm{At}(P,A)$ are both in $I$; the negative
precondition $\mathrm{Holding}(P)$ is absent, and absent means false.

\item \textbf{Is $\mathrm{Drop}(Package,C)$ applicable in $I$?}
No, and it fails on \emph{both} preconditions --- the robot is not at C and is
not holding anything. It is tempting to say yes because it is the action that
achieves the goal and it is on the list of available actions. Neither of those
is the test.

\item \textbf{What is the missing box in Task 4's diagram?}
The \textbf{applicability decision}: is $S \ent \mathrm{Pre}(a)$? Everything
above it is logic; everything below it is search.

\item \textbf{Where is logical reasoning used?}
Four places: the applicability test $S \ent \mathrm{Pre}(a)$; the state update
$S' = (S \setminus \mathrm{Del}) \cup \mathrm{Add}$; the goal test $S \ent G$;
and --- outside the agent entirely --- the Prolog verifier, where logic judges
the agent's output rather than driving it.

\item \textbf{An error from not checking preconditions?}
The planner accepts $\mathrm{Move}(A,B)$ then $\mathrm{PickUp}(Package,B)$ and
``picks up'' a package that is not there. The resulting state contains
$\mathrm{Holding}(P)$ \emph{and} $\mathrm{At}(P,A)$ --- the package is in two
places at once, and every later state is nonsense.

\item \textbf{Why is planning the same as the search module?}
Same frontier, same visited set, same parent map, same ``BFS is optimal for
unit costs''. The only difference is that the successor function is
$\{\mathrm{Apply}(S,a) : a \in A,\ S \ent \mathrm{Pre}(a)\}$, computed by a
logical test, instead of being read off a map. Two consequences: the state space
is implicit and exponential (so real planners derive heuristics from the action
descriptions --- relaxed-plan heuristics ignore delete lists --- rather than
being handed an $h$ like A$^\ast$ was); and the \emph{shape of the graph} is
now something you wrote, and therefore something you can get wrong.

\item \textbf{What did the LLM contribute, and what did you verify?}
It contributed the engineering: the class structure, the one-line set operations
for applicability and effects, the BFS driver, the printer, and correct
diagnoses once symptoms were described. Verified independently: that the plan
re-executes from $I$; that the planner terminates on an unsolvable problem; that
the plan is shortest (against a separately written iterative-deepening oracle);
and --- the one that could not be done in Python at all --- that the action
model matches the warehouse.

\end{enumerate}

% =====================================================================
\section{The five words to remember}

\begin{center}
\Large
Understand $\rightarrow$ Specify $\rightarrow$ Generate $\rightarrow$ Execute $\rightarrow$ Verify
\end{center}

\vspace{4pt}

\begin{center}
\small
\begin{tabular}{@{}ll@{}}
\toprule
\textbf{Stage} & \textbf{What it caught in this laboratory} \\
\midrule
Understand & that the goal is partial, so the goal test is $\subseteq$ \\
Specify    & Prompt 1 made Defect 0 a one-line fix rather than a puzzle \\
Generate   & a correct BFS planner, quickly \\
Execute    & Test B found Defect 1 --- a hang, invisible to Test A \\
Verify     & Prolog found Defect 2 --- invisible to \emph{everything} inside Python \\
\bottomrule
\end{tabular}
\end{center}

\vspace{8pt}

\begin{keyidea}{The last word}
An AI system can generate a candidate solution, while a separate logical system
checks it.

The check has to come from \textbf{outside} the thing being checked. A program
that grades its own homework will pass --- and so will its explanation of why it
deserved to.
\end{keyidea}

\end{document}
